% Source: https://artofproblemsolving.com/community/c3h2514313_question_about_fractional_exponents
% Stable identity: aops-topic-2514313-q01
\textbf{(a)} By the definitions,
\[4^{\dfrac12}=\sqrt4=2,\qquad\sqrt[4]{16}=2,\qquad8^{\dfrac23}=\left(\sqrt[3]{8}\right)^2=2^2=4,\qquad16^{-\dfrac34}=\dfrac{1}{\left(\sqrt[4]{16}\right)^3}=\dfrac{1}{8}.\]
In particular $4^{\dfrac12}=2$, not $\pm2$: like $\sqrt4$, the power $4^{\dfrac12}$ names one number, the positive square root.

\textbf{(b)} Factor:
\[x^4-16=\left(x^2-4\right)\left(x^2+4\right)=(x-2)(x+2)\left(x^2+4\right).\]
Since $x^2+4>0$ for real $x$, the real solutions are $x=2$ and $x=-2$.

The symbol $\sqrt[4]{16}$ is defined to be one particular number (the positive one), so that expressions such as $\sqrt[4]{16}+1$ have a single value. The equation $x^4=16$ asks instead for \emph{every} number whose fourth power is $16$, and there are two: $\sqrt[4]{16}$ and $-\sqrt[4]{16}$. In the same way $\sqrt4=2$, while $x^2=4$ has the two solutions $\pm2$.

\[\boxed{\text{(a) }2,\ 2,\ 4,\ \dfrac18;\qquad\text{(b) }x=2\text{ or }x=-2}\]
